\chapter{Conclusion}
Nowadays, flight testing is becoming crucial in studying aeroelastic phenomena, due to limitations of wind tunnels. However, existing flying laboratories for aeroelastic research are often limited in versatility and financially infeasible for many universities and research laboratories. Therefore, the purpose of this report was to present the design of the High-speed Unmanned General-purpose Observatory (HUGO), a reusable test platform for high-speed aeroelastic research that allows for testing multiple wing planforms and aircraft configurations.

The market analysis identified research centres such as NLR in the Netherlands and DLR in Germany, and universities such as TU Delft as key customers. The cost of using HUGO was estimates to €25,000 per flight day excluding setup costs. Consequently, it was estimated that this aircraft would take 20\% of the total market. Due to the needs of the target customers, Den Helder, North Sea, controlled airspace), technical risks analysis and mitigation, Type Certificate (classification under SC-Light-UAS High risk, but flight envelope already compliant with CS-23 aerobatic - more constraining), Design Organisation Approval for not having EASA to certify each and every planform/configuration and classifying major changes as research procedures,

\textcolor{red}{Writing in progress...}
The starting point was the aircraft configuration selected in the previous phase of the project. Namely, a high wing aircraft featuring static ports for wing and canard attachment, removable T-tail cone, buried retractable landing gear and fuselage-mounted engines.




The envelope of wing configurations (\textcolor{red}{to be defined here!}) defined by span, sweep angle, taper. The structure, internal and external layout are all design to accommodate all of those planforms.
Also the standard planform defined (free of flutter, NASA SC(2)-0012 supercritical airfoil)

uses SAF, reaches Mach 0.75 for 5 minutes at FL270, collects measurement data, stable and unstable configurations

 sustainability(contribution to future load alleviation systems, effectively to reduction of emissions, SAF) Tests: AFS, GLA, aileron reversal (external control system), folding wingtips (external wing), truss-braced wing ...

Technical design
external/internal layout. 
Manufacturing: metals, composites, core material chosen based on mechanical performance, environmental resistance, low-volume aerospace-grade manufacturability.
The primary structure utilizes out-of-autoclave (OOA) CYCOM
5320-1 carbon fiber epoxy prepreg, fuselage honeycomb structure. Fuselage assembly: lap joints between inner and outer skins, adhesive film for the core, pylon-to-fusleage connection characterised by mechanical fastening of cylindrical metal (titanium) inserts, landig gear assembly through mechanical fastening Structure sized for ultimate load factor of 9g

Structural analysis and design: wing stresses, buckling and deflections based on lift distribution, fuselage under pull-up manoeuvre and 2-point landing condition, abuse loads, aeroelastic analysis encompassing divergence and flutter. Wing thickness of 2.5mm as the minimum, fuselage sheet thickness of ...mm.
\textcolor{red}{1-2 sentences about noise and environmental impact}
Last but not least, emennage and canard have been sized for neutral static stability to enable testing of stable, neutrally stable and unstable flight configurations. Regarding the dynamic stability, the aircraft equipped with the standard planform exhibited acceptable stability characteristics.

Standard payload: dual LIDAR sensors for in-flight deflection measurements, strain gauges, cameras. WIng-fuselage connection through four 6-axis load cells, space for additional sensors etc., two JetCat P100-RX-BL micro-turbine engines
dynamically shifting the CG through fuel management system, high-performance steel for the landing gear,

Finally, ...